\documentclass{tietreport}

% Bibliography configuration
\usepackage[
  style=alphabetic,
  backend=biber,
  sorting=nty,
  maxnames=5,
  minnames=3,
  backref=true,
  doi=false,
  isbn=false,
  url=false,
  eprint=false
]{biblatex}
\addbibresource{references.bib}

%% ==================================================
%% CUSTOMIZE THESE FIELDS FOR YOUR PROJECT
%% ==================================================

% Course/Lab header
\TitlePageHeader{%
  Your Course Code - Course Title%
}

% Main project title
\ProjectTitle{%
  Your Project Title Here%
}

% Subtitle or project description
\ProjectSubTitle{%
  A Descriptive Subtitle for Your Work%
}

% Student information (up to 4 students)
\author{%
  \texttt{102XXXXXXX} Student Name One \\
  \texttt{102XXXXXXX} Student Name Two \\
  \texttt{102XXXXXXX} Student Name Three \\
  \texttt{102XXXXXXX} Student Name Four%
}

% Group number, degree, year, program
\TitlePageSubText{%
  Group: \texttt{XXXX} BE Third Year, CSE%
}

% Faculty advisor/guide name
\AdvisorName{%
  Dr. Faculty Member Name%
}

% Evaluation period and department info
\TitlePageFooterText{{%
    \large Mid-Semester Evaluation \\[0.2cm]
    \bfseries Your Department Name%
  } \\[0.15cm]
  Thapar Institute of Engineering and
  Technology, Patiala%
}

%% ==================================================

\begin{document}

% Generate title page
\maketitle

% Table of contents (auto-generated from chapters)
\tableofcontents

% ===== CHAPTER 1: INTRODUCTION =====

\chapter{Introduction}

\section{Background and Context}

Begin your introduction here. Provide context
for your project and explain why it matters.
This section should clearly establish the
problem or opportunity your project addresses.

\subsection{Problem Statement}

Define the specific problem your project
tackles. Be precise and provide relevant
background information that helps readers
understand the scope and significance.

\section{Project Objectives}

\begin{enumerate}
  \item \textbf{Objective 1:} State your first
    specific, measurable objective.
  \item \textbf{Objective 2:} State your second
    objective.
  \item \textbf{Objective 3:} State your third
    objective.
\end{enumerate}

Document your SMART objectives (Specific,
Measurable, Achievable, Relevant, Time-bound)
clearly for your evaluators.

% ===== CHAPTER 2: METHODOLOGY & DESIGN =====

\chapter{Methodology and Design}

\section{System Architecture}

Describe your overall system design. You can
reference figures and tables in your document:

\subsection{High-Level Design}

Explain the major components and how they
interact. This section demonstrates your
understanding of the problem and your
solution approach.

\subsection{Technical Stack}

\begin{itemize}
  \item \textbf{Framework/Language:} List your
    primary technologies.
  \item \textbf{Database:} Specify database
    systems if applicable.
  \item \textbf{Tools:} List development and
    testing tools used.
  \item \textbf{Libraries:} Note key libraries
    or frameworks.
\end{itemize}

\section{Implementation Details}

\subsection{Module 1: Core Functionality}

Describe the implementation of your first
major module. Include technical details about
algorithms, data structures, or design
patterns used.

\subsection{Module 2: Secondary Features}

Continue documenting your implementation
by module, demonstrating the complexity and
depth of your work.

% ===== CHAPTER 3: RESULTS & EVALUATION =====

\chapter{Results and Evaluation}

\section{Testing Methodology}

Describe your testing strategy:

\begin{itemize}
  \item \textbf{Unit Testing:} How you tested
    individual components.
  \item \textbf{Integration Testing:} How you
    verified components work together.
  \item \textbf{Performance Testing:} How you
    measured system performance.
  \item \textbf{User Testing:} If applicable,
    how users evaluated your system.
\end{itemize}

\section{Performance Metrics}

Present your results clearly with tables and
figures. For example:

\begin{table}[h]
\centering
\begin{tabular}{|l|r|r|}
  \hline
  \textbf{Metric} & \textbf{Target} &
  \textbf{Achieved} \\
  \hline
  Response Time & $< 100$ms & $85$ms \\
  \hline
  Accuracy & $> 95\%$ & $97.2\%$ \\
  \hline
  Throughput & $1000$ req/s & $1250$ req/s \\
  \hline
\end{tabular}
\caption{Performance Results Summary}
\label{tab:performance}
\end{table}

\section{Analysis}

Analyze your results:

\begin{itemize}
  \item Did you meet your objectives?
  \item Where did you exceed expectations?
  \item What challenges did you encounter?
  \item How did you resolve them?
\end{itemize}

% ===== CHAPTER 4: CONCLUSION =====

\chapter{Conclusion}

\section{Summary of Work}

Summarize what you accomplished. Remind
readers of your objectives and explain which
ones you met.

\section{Key Findings}

\begin{itemize}
  \item \textbf{Finding 1:} State a significant
    result or insight from your work.
  \item \textbf{Finding 2:} Describe another
    important discovery.
  \item \textbf{Finding 3:} Highlight any
    unexpected results.
\end{itemize}

\section{Future Work and Recommendations}

Discuss how others could extend or improve
your work:

\begin{enumerate}
  \item Potential enhancements to your system
  \item Alternative approaches worth exploring
  \item Scalability improvements
  \item Integration with other systems
  \item Real-world deployment considerations
\end{enumerate}

\section{Final Remarks}

Conclude with reflections on what you learned
and the impact of your work.

% ===== BIBLIOGRAPHY =====

% Print bibliography with heading in TOC
\printbibliography[heading=bibintoc]

\end{document}
